\chapter{La química del vínculo}

\epigraph{El enamoramiento es la forma que tiene la naturaleza de que dos personas se miren el tiempo suficiente como para decidir quedarse.}{Helen Fisher}

Casi todo el mundo reconoce la experiencia, aunque no sepa explicarla. Te ha pasado, o te está pasando, o conoces a alguien a quien le pasa: de repente, una persona que antes era una más entre miles ocupa el centro de todo. No puedes dejar de pensar en ella. La recuerdas al despertar y al acostarte. Cada mensaje suyo te produce un pequeño estallido de alegría. Comes menos, duermes peor, y sin embargo te sientes lleno de energía. El mundo entero parece más luminoso, más intenso, más lleno de sentido.

Estás enamorado.

Y lo que sientes no es una metáfora. Es química. Literalmente.

\section{El cerebro enamorado}

En el año 2005, la antropóloga Helen Fisher y su equipo publicaron un estudio que marcó un antes y un después en la comprensión científica del amor.\footnote{Fisher, H., Aron, A., y Brown, L.L., ``Romantic love: a mammalian brain system for mate choice'', \textit{Philosophical Transactions of the Royal Society B}, 2006. Fisher llevaba décadas estudiando el amor desde la antropología y la neurociencia, y es autora de varios libros sobre el tema, entre ellos \textit{Why We Love} (2004) y \textit{Anatomy of Love} (1992).} Tomaron a un grupo de personas que declaraban estar intensamente enamoradas, las metieron en un escáner de resonancia magnética funcional, y les mostraron fotografías de la persona amada. Lo que vieron en las pantallas fue revelador.

El cerebro de una persona enamorada se ilumina de una forma muy específica. No se activan principalmente las zonas de la emoción, como cabría esperar. Lo que se enciende con fuerza es el \textit{área tegmental ventral} y el \textit{núcleo caudado}: regiones profundas del cerebro, muy antiguas evolutivamente, que forman parte del llamado sistema de recompensa.\footnote{El sistema de recompensa es un circuito cerebral que refuerza las conductas necesarias para la supervivencia ---comer, beber, reproducirse--- generando una sensación de placer que nos impulsa a repetirlas. Que el enamoramiento active este sistema sugiere que, para el cerebro, el vínculo de pareja es tan esencial como alimentarse.} Son las mismas zonas que se activan ante una recompensa intensa, ante algo que el cerebro interpreta como vital, como necesario. Como algo que no puede dejar pasar.

Y el combustible de esa activación tiene nombre: \textbf{dopamina}.

\section{El cóctel de la pasión}

La dopamina es un neurotransmisor ---una molécula que las neuronas usan para comunicarse entre sí--- asociado a la motivación, la búsqueda de recompensa y el deseo. Cuando estamos enamorados, el cerebro libera dopamina en cantidades muy superiores a lo habitual. Es la dopamina la que produce esa sensación de euforia, esa energía inagotable, ese enfoque casi obsesivo en la persona amada.\footnote{Los niveles de dopamina en el enamoramiento se han comparado con los que producen ciertas drogas estimulantes. No es una exageración: el sistema de recompensa activado es el mismo. La diferencia es que el enamoramiento activa el circuito de forma natural y con un propósito biológico. Cf. Fisher, H., \textit{Why We Love: The Nature and Chemistry of Romantic Love}, Henry Holt, 2004.}

Pero la dopamina no actúa sola. Junto a ella sube la \textbf{noradrenalina}, que pone al cuerpo en estado de alerta ---de ahí las manos sudorosas, el corazón acelerado, la dificultad para dormir, la pérdida de apetito---. Y ocurre algo más, quizá lo más revelador: baja la \textbf{serotonina}.

La serotonina es el neurotransmisor de la calma, de la regulación emocional, del equilibrio. Cuando sus niveles descienden, aparecen los pensamientos intrusivos, la rumiación, la incapacidad de dejar de pensar en algo. ¿Les suena? Un estudio célebre de la psiquiatra Donatella Marazziti, de la Universidad de Pisa, comparó los niveles de serotonina de personas recién enamoradas con los de pacientes diagnosticados de trastorno obsesivo-compulsivo. El resultado fue asombroso: eran prácticamente iguales.\footnote{Marazziti, D., Akiskal, H.S., Rossi, A., y Cassano, G.B., ``Alteration of the platelet serotonin transporter in romantic love'', \textit{Psychological Medicine}, 1999. Los enamorados presentaban niveles de serotonina un 40\% más bajos que los del grupo de control, similares a los de pacientes con TOC.}

Sí: el enamoramiento, bioquímicamente hablando, se parece a una obsesión.

No dejas de pensar en esa persona. La idealizas. Repasas mentalmente cada conversación, cada gesto, cada palabra. Cualquier estímulo ---una canción, un olor, una calle--- te la trae a la memoria. Tu capacidad de concentrarte en otras cosas disminuye notablemente. Si alguien te dijera que te comportas como un obsesivo, te parecería exagerado. Pero tu cerebro, si pudiera hablar, confirmaría el diagnóstico.

\section{Una locura con sentido}

Ahora bien: antes de alarmarnos, conviene hacer una pregunta. ¿Por qué? ¿Por qué la evolución diseñó un estado tan intenso, tan absorbente, tan parecido a la locura?

La respuesta tiene que ver con algo que ya vimos en el primer capítulo: nuestros bebés nacen extraordinariamente vulnerables. Necesitan años de cuidado. Y para que una cría humana sobreviva, no basta con una madre: hace falta una pareja que se comprometa. Que se quede. Que colabore.

El enamoramiento es el mecanismo que la naturaleza utiliza para que dos personas se elijan mutuamente y se mantengan juntas el tiempo suficiente como para establecer un vínculo profundo. Es una especie de turbo biológico: durante unos meses ---a veces algo más de un año, raramente más de tres--- el cerebro concentra toda su maquinaria motivacional en una sola persona. Te hace sentir que esa persona es única, insustituible, la más importante del mundo. Y en cierto modo tiene razón: porque lo que está en juego, evolutivamente, es la supervivencia de la especie.

Helen Fisher propuso que el amor humano funciona en realidad como tres sistemas cerebrales interconectados pero distintos:\footnote{Fisher, H., ``Lust, Attraction, and Attachment in Mammalian Reproduction'', \textit{Human Nature}, 1998. Fisher desarrolló ampliamente este modelo en \textit{Why We Love} (2004) y \textit{Anatomy of Love} (edición revisada, 2016).}

\begin{itemize}
    \item El \textbf{deseo sexual} (\textit{lust}), impulsado fundamentalmente por las hormonas sexuales ---testosterona y estrógenos---, que nos orienta hacia la búsqueda de pareja en general.
    \item La \textbf{atracción romántica} (\textit{attraction}), mediada por la dopamina y la noradrenalina, que concentra nuestra atención en una persona concreta y genera el estado que llamamos enamoramiento.
    \item El \textbf{apego de pareja} (\textit{attachment}), sostenido por la oxitocina y la vasopresina, que produce la sensación de calma, seguridad y conexión profunda con esa persona a largo plazo.
\end{itemize}

Los tres sistemas pueden activarse juntos, pero no siempre lo hacen. Puedes sentir deseo por alguien sin estar enamorado. Puedes estar enamorado de alguien y sentir deseo por otro. Puedes sentir apego profundo por tu pareja de muchos años sin experimentar ya la tormenta del enamoramiento inicial. La vida sentimental humana es compleja precisamente porque estos tres sistemas no siempre van de la mano.

Pero cuando los tres coinciden en una misma persona ---cuando deseas, te enamoras y te vinculas con la misma persona---, algo extraordinario ocurre. Y es ahí donde entra la segunda parte de la historia.

\section{De la tormenta a la calma}

El enamoramiento no dura para siempre. No puede durar para siempre. Un estado de activación tan intenso consume una cantidad enorme de recursos. Si el cerebro mantuviera indefinidamente los niveles de dopamina del enamoramiento, acabaríamos agotados, incapaces de concentrarnos en otra cosa, funcionando al límite.\footnote{Diversos estudios estiman que la fase de enamoramiento intenso dura entre 12 y 18 meses, aunque con variaciones individuales importantes. Cf. Tennov, D., \textit{Love and Limerence: The Experience of Being in Love}, Stein and Day, 1979.}

Y aquí es donde muchas parejas se asustan. Porque cuando la tormenta de dopamina empieza a amainar, cuando ya puedes dormir con normalidad y recuperas el apetito, cuando ya no sientes ese vuelco en el estómago cada vez que ves a tu pareja, es tentador pensar que el amor se ha terminado. Que ``se ha apagado la llama''. Que ya no sientes lo mismo.

Pero lo que ocurre es algo muy distinto. Lo que ocurre es que el cerebro está haciendo una transición. Está pasando del sistema de atracción al sistema de apego. Y las protagonistas ahora son otras moléculas: la \textbf{oxitocina} y la \textbf{vasopresina}.

Ya encontramos la oxitocina en el capítulo anterior, vinculada al parto y la lactancia. Pero su papel va mucho más allá. La oxitocina se libera con el contacto físico ---un abrazo, una caricia, la relación sexual---, y genera sensaciones de confianza, calma y bienestar.\footnote{Cf. Uvnäs-Moberg, K., \textit{The Oxytocin Factor: Tapping the Hormone of Calm, Love, and Healing}, Da Capo Press, 2003.} No produce el estallido eufórico de la dopamina. Produce algo más silencioso pero más estable: la sensación de estar en casa. De estar con la persona adecuada. De seguridad.

La vasopresina, por su parte, está especialmente vinculada a la conducta protectora y al vínculo de pareja en los varones. Uno de los estudios más fascinantes sobre esta molécula se realizó con un animal inesperado: el topillo de las praderas.\footnote{Young, L.J. y Wang, Z., ``The neurobiology of pair bonding'', \textit{Nature Neuroscience}, 2004. Los topillos de las praderas (\textit{Microtus ochrogaster}) son uno de los pocos mamíferos monógamos. Sus primos, los topillos de montaña (\textit{Microtus montanus}), son promiscuos. La diferencia clave está en la distribución de receptores de vasopresina en el cerebro.} Los topillos de las praderas son uno de los pocos mamíferos que forman parejas estables de por vida. ¿La clave? Tienen una alta densidad de receptores de vasopresina en las zonas de recompensa del cerebro. Cuando se les bloquean esos receptores, dejan de formar vínculos de pareja. Cuando se aumentan artificialmente en una especie de topillo promiscuo, este empieza a vincularse con una sola hembra.

No somos topillos, por supuesto. Pero la molécula es la misma, y el mecanismo básico se conserva: la vasopresina facilita la fidelidad y la protección del vínculo.

\section{El amor que permanece}

Lo que la ciencia nos muestra, entonces, es que el amor humano no es un único fenómeno, sino un proceso. Un proceso que empieza con una tormenta y, si todo va bien, desemboca en un puerto.

La tormenta es espectacular. Es intensa, es absorbente, es memorable. Pero es, por definición, temporal. El puerto no tiene la misma adrenalina, pero tiene algo que la tormenta no puede dar: profundidad, estabilidad, la certeza de que hay alguien a tu lado cuando la vida se pone difícil.

Y aquí viene lo verdaderamente interesante: el apego no es una versión rebajada del enamoramiento. No es ``conformarse'' con menos. Es un sistema biológico distinto, diseñado para un propósito distinto. El enamoramiento sirve para elegir; el apego sirve para permanecer. Y ambos son necesarios.

Los estudios de parejas que llevan décadas juntas y que declaran seguir profundamente enamoradas muestran algo notable: en sus cerebros, los circuitos de apego están muy activos, como cabría esperar, pero además siguen mostrando activación en las áreas de recompensa asociadas a la dopamina.\footnote{Acevedo, B.P., Aron, A., Fisher, H.E., y Brown, L.L., ``Neural correlates of long-term intense romantic love'', \textit{Social Cognitive and Affective Neuroscience}, 2012. Este estudio escaneó los cerebros de personas que llevaban una media de 21 años de relación y seguían reportando amor intenso. Las áreas de recompensa seguían activas, pero sin la activación de las áreas de ansiedad y obsesión típicas del enamoramiento temprano.} Es decir: el amor maduro conserva la intensidad, pero pierde la ansiedad. Se parece más a una llama constante que a un incendio.

\section{Más allá de la química}

Sería tentador detenernos aquí y concluir que el amor es, al final, solo una cuestión de moléculas. Que somos marionetas de la dopamina y la oxitocina, y que lo que llamamos amor no es más que un sofisticado programa biológico de vinculación.

Pero hay algo que no encaja en esa explicación. Algo que la biología sola no puede explicar.

Porque la biología nos empuja, sí. Nos pone el viento a favor. Pero no nos obliga. Hay personas que sienten el torbellino del enamoramiento y deciden no seguirlo. Hay parejas que pierden el apego biológico y eligen quedarse. Hay hombres y mujeres que, contra toda la química de su cuerpo, deciden ser fieles a una promesa que hicieron un día. Y hay otros que, con toda la oxitocina del mundo circulando por sus venas, deciden marcharse.

El ser humano es el único animal que puede decirle que no a sus hormonas. El único que puede elegir contra su instinto. Y el único que puede elegir \textit{a favor} de su instinto no porque le arrastre, sino porque libremente decide que esa es la dirección correcta.

Eso es lo que hace al amor humano algo radicalmente distinto al vínculo de pareja de los topillos. No es que nuestra oxitocina sea diferente. Es que nosotros somos más que nuestra oxitocina. Somos capaces de tomar esa corriente biológica que nos empuja hacia el otro y convertirla en una decisión. En un proyecto. En una promesa.

El enamoramiento te dice: ``No puedo dejar de pensar en ti''. El apego te dice: ``Me siento seguro a tu lado''. Pero hay algo más ---algo que no se reduce a la bioquímica--- que dice: ``Te elijo. Hoy. Y mañana. Y pasado mañana. No porque mis hormonas me lo pidan, sino porque quiero construir mi vida contigo''.

Ese ``algo más'' es, quizá, lo más humano que hay en nosotros. La biología nos regala un impulso poderoso. Pero convertir ese impulso en un amor que dura toda la vida requiere algo que ninguna molécula puede proporcionar: la libertad de elegir.

Y aquí se abre una pregunta que merece la pena hacerse. Si nuestro cuerpo está tan finamente diseñado para el vínculo ---con sistemas que nos empujan a elegir, a apegarnos, a permanecer---, y si al mismo tiempo somos libres para acoger o rechazar ese impulso... ¿de dónde viene esa doble dotación? ¿Quién nos hizo capaces de amar \textit{y} de elegir amar?

Pero eso es materia para más adelante. De momento, basta con asombrarse.

\paracaminar

\begin{enumerate}
  \item Recordad juntos los primeros tiempos de vuestra relación: las mariposas, la dificultad para pensar en otra cosa, la euforia. Ahora que sabéis que había dopamina, noradrenalina y serotonina baja detrás de todo aquello, ¿cambia en algo lo que sentisteis? ¿O lo hace más fascinante?

  \item ¿Habéis experimentado ya la transición del enamoramiento al apego? Si es así, ¿cómo la vivisteis? ¿Hubo un momento en que pensasteis que el amor se estaba ``apagando'', o supisteis reconocer que estaba cambiando de forma?

  \item Hablad juntos: ¿qué significa para vosotros pasar de ``no puedo dejar de pensar en ti'' a ``te elijo cada día''? ¿Qué tiene de distinto un amor elegido frente a un amor sentido?
\end{enumerate}

\vinetacapitulo{ch03}
